% ============================================================
%  P2 paper skeleton v0.2  ·  BioSNN-Plug Benchmark Paper
%  Build:  pdflatex main && bibtex main && pdflatex main && pdflatex main
%         (or simply:  latexmk -pdf main)
%  写作纪律: 数字全部来自 research/ 复现记录; \todo{} 标记待填
%  红线: 不得出现 "state-of-the-art"; 预训练组件必须声明冻结
% ============================================================
\documentclass[11pt]{article}
\usepackage[margin=1in]{geometry}
\usepackage[T1]{fontenc}
\usepackage[utf8]{inputenc}
\usepackage{natbib}   % \citet/\citep；须在 hyperref 之前载入。plainnat.bst 随 TinyTeX 发布
\usepackage{booktabs, graphicx, hyperref, amsmath, xcolor}
\newcommand{\todo}[1]{\textcolor{red}{[TODO: #1]}}
\title{Local Learning Rules Under Honest Measurement:\\A Reproducibility-First Benchmark of e-prop, Hebbian, and R-STDP\\on Consumer Hardware}
\author{lzmd-arch\thanks{Independent researcher. Correspondence: \todo{email}.}\\ \small{github.com/lzmd-arch/BioSNN-Plug}}
\date{October 2026}

\begin{document}
\maketitle

% ------------------------------------------------------------
\begin{abstract}
\todo{120-180 words. Outline: (1) problem -- local-rule benchmarks lack honest measurement (budget/protocol/degenerate solutions); (2) what we did -- three rule lines + interaction matrix + same-budget BPTT control, consumer GPU; (3) headline numbers -- e-prop ALIF converges at 0.8831 (250ep); LIF-vs-ALIF interaction; Critic-corrected R-STDP 237.8 steps; (4) methodological contributions -- bit-exact reproduction, retention-joint reporting, alpha-to-zero degeneracy analysis; (5) stance -- gaps reported as-is, no cosmetics.}
\end{abstract}

% ------------------------------------------------------------
\section{Introduction}
\subsection{Motivation}          % pain point: BPTT biological implausibility + memory/compute-graph cost
\subsection{The measurement gap} % status quo: local-rule papers use incomparable budgets/protocols
\subsection{Contributions}
\begin{itemize}
  \item \textbf{Honest benchmark}: three local-rule lines under one unified protocol (budget/seeds/protocols fully public). \todo{fill numbers}
  \item \textbf{Interaction finding}: neuron type x learning rate swept as a controlled grid; the adaptation term is a net drag at every learning rate tested (same-budget LIF leads ALIF by 0.93-1.23pp over lr 1e-3 to 1e-2; 5 seeds each).
  \item \textbf{Convergence evidence}: e-prop + ALIF converges at 0.8831 (250 epochs), lower edge of the e-prop literature band.
  \item \textbf{Methodology}: bit-exact reproduction, retention-joint reporting, degeneracy analysis, frozen-noise evaluation protocol.
  \item \textbf{Open assets}: all code/configs/scripts open (biosnn-bus on PyPI); comparison table with evidence ledger for 5 sibling projects.
\end{itemize}
\subsection{Positioning: what this paper is not}
% Three non-claims: no SOTA claims; no emergence/intelligence claims; no hidden negative results (SHD reported as-is).

% ------------------------------------------------------------
\section{Related Work}
\subsection{Local learning rules for spiking networks}
A family of learning rules for spiking neural networks (SNNs) restricts synaptic
updates to quantities available at the synapse --- presynaptic and postsynaptic
activity together with a broadcast modulation signal --- thereby avoiding
backpropagation through time (BPTT) and its requirement to store complete
activation histories. We refer to these as \emph{local learning rules} and
review the five most directly relevant lines of work.

\citet{bellec2020eprop} introduced e-prop, a three-factor rule in which each
synapse maintains a decaying eligibility trace computed from local spike
events, while a learning signal broadcast from the readout layer supplies the
third factor. The rule requires O(1) memory in sequence length and was shown to
approach BPTT-level performance on speech and sequence tasks with adaptive
LIF (ALIF) neurons. Published e-prop results, however, are typically reported
at a single hyperparameter setting with a small number of seeds, leaving the
rule's sensitivity to optimization choices largely uncharacterized. In this
work we treat such sensitivity as a first-class measurement target rather
than a nuisance variable.

\citet{pes2025tp} move the trace from the synapse to the neuron: traces
propagation (TP) keeps two decaying activity traces per neuron rather than an
eligibility trace per synapse, and derives the learning signal from a layer-wise
contrastive loss, for a space complexity of $O(LH)$ against the $LH^{2}$ the
same paper tabulates for e-prop. Its measurements, however, are taken on fully
connected feedforward architectures with LIF neurons, and part of the SHD
comparison it prints is transcribed from a third party rather than re-run under
its own protocol --- the paper footnotes those rows as such. TP therefore does
not establish that the per-synapse traces which make e-prop's storage quadratic
can be avoided in a recurrent core, which is the setting we measure; on the
paper's own N-MNIST table it is e-prop, not TP, that takes the higher number
(97.90 against 97.33). We accordingly treat TP as a parallel rule rather than a
drop-in replacement, and do not use it as the storage argument for our
recurrent model.

\citet{espp2024} take a different route: EchoSpike predictive plasticity (ESPP)
uses the whole spike train of the previous sample as a prediction target, and
pulls same-label samples together while pushing different-label ones apart, a
layer-local rule that never invokes automatic differentiation. The rule also
decides for itself which time steps are worth updating --- a threshold on input
activity plus a threshold on the loss --- and the authors report that weight
updates are executed on only 18--27\% of time steps, a fraction that decreases
during training, alongside an SHD accuracy of 84.32\%. Both figures come from a
single configuration and are self-reported; we are not aware of an independent
reproduction. ESPP is also not an e-prop variant: it is a
predictive/contrastive layer-local rule, and the selective updating it performs
is a different mechanism from the event-driven form discussed next, with which
it is easily confused.

\citet{korcsak2025event} keep the rule and change when it is evaluated: a
synapse updates only when it receives a spike event, and the implementation is
embedded in a large-scale simulation framework (NEST). It reproduces the
original rule's performance on pattern generation, evidence accumulation and
N-MNIST, and weak- and strong-scales to two million neurons on sparse networks;
to satisfy strict locality it also removes one violation of that principle in
the original rule --- the eligibility filter as published depends on the time
constant of the output neuron, which a synapse cannot know. Those are
measurements taken at HPC scale, and the regime we can afford is a single
consumer GPU with networks roughly four orders of magnitude smaller, where the
question is not whether event-driven updating scales but whether it changes the
numbers at all. Closer to our setting, this is the variant we would reach for
first if scale became the binding constraint; it is not, however, the object of
the benchmark below, which covers the form we can reproduce bit-for-bit on one
machine.

\citet{millidge2025deep} generalize the trace recursion itself, from a single
recurrent layer to arbitrary depth and arbitrary DAGs, arguing that the cost
stays linear in depth. The note contains no experiments --- its author states
that none were run --- so it can support the methodological claim that a
multi-layer recurrent SNN is trainable by purely local rules in principle, and
cannot support any accuracy figure. Its own stated limitations cut against the
word ``online'': e-prop as specified updates weights at the end of an episode
rather than during it, and keeping one set of traces per parameter group becomes,
in the author's phrase, quickly unmanageable with depth. We cite it for those two
statements and for nothing measured.

\subsection{Meta-learned plasticity and arbitration}
A second line does not ask which local rule to use but how the rules themselves
should change: plasticity becomes a learned or slowly varying quantity.

\citet{confavreux2025} meta-learn co-activity plasticity rules for recurrent
spiking networks with an evolutionary-strategy outer loop, over
parameterizations of increasing complexity. Their result is itself a caution:
beyond a certain complexity the meta-learned strategies begin to fail, and
specifying a loss that constrains the emerging dynamics a priori is hard. We do
not meta-learn --- the rules stay fixed and any arbitration between them is
specified by hand --- so the question we can measure is a different one: what a
fixed mixture buys at a matched optimization budget.

\citet{miconi2018differentiable} sit at the other end of the same axis, training
the plasticity coefficients themselves by backpropagation through a plastic
network. That is meta-learning with a global gradient, and we cite it as a
contrast rather than as support: its results say nothing about what purely local
rules achieve, which is the question this paper is built around.

\citet{zanardi2024metaplasticity} analyse metaplasticity as a memory mechanism
rather than as a training algorithm: on a cyclic feed-forward network with
several levels, where a Hebbian term and a metaplastic reinforcement term act on
grouped edges, the dynamics separate into three regimes, and in two of them a
stored path can still be retrieved after the weights have been reset to zero ---
retention that survives at the slower level after the fast one has been erased.
That is the mechanism behind the retention constraint we impose when rules are
mixed; it supplies no accuracy number, and we borrow none.

\citet{turcu2025switching} are the closest neighbour. Their synapses change
their governing plasticity rule independently during training, switching between
a weak Hebbian-like pre-post rule for weak synapses and a strong
credit-assignment rule (backpropagation) for strong ones; the switching network
is reported to learn with fewer examples than a backpropagation-only baseline,
and to develop a more antisymmetric recurrent weight matrix. Two differences
separate it from our setting. They switch discretely between rules where we mix
two rule updates continuously, and their strong rule is backpropagation, so the
mixture is not itself purely local. Closer to our setting, the question that
remains is the one we pose: what must arbitration satisfy so that mixing two
updates does not degrade what either rule has already retained.

\subsection{SNN--LLM hybrid architectures}               % EMBER; MEMBRAIN; Javis (quote its learning-rule status verbatim and honestly)
  % 未写：MEMBRAIN [10] 与 Javis [16] 是 docs/references.md 里的条目，但 refs.bib 只收了
  % EMBER；按 refs.bib 顶部的规则，这三个键要一起动才能写这一段。
\subsection{Reproducibility protocols in SNN training}  % sparse literature; our Section 3 fills the gap
% One sentence per citation: "difference from ours". Grouping lives in docs/references.md
% and docs/related_projects.md.（原注释指向 docs/related_work.md，该文件在本仓不存在。）
  % 未写：D-RTRL / pp-prop —— refs.bib 无对应条目，且按顶部规则不得新增仓库外引用。

% ------------------------------------------------------------
\section{Methods}
\subsection{Systems under test}   % W1 Hebbian (kernelized + divisive norm.); W2 e-prop (ALIF/LIF); W3 R-STDP + TD-LTP Critic
\subsection{Evaluation protocol}  % this section is a selling point -- write it thick:
  % (a) budget discipline: all comparisons at matched epochs/sample budgets; cross-budget numbers must never be subtracted;
  % (b) config dumps: mandatory per-run machine-readable config; every paper number is queryable;
  % (c) seeds & bit-exact reproduction: 5 seeds; identical results across sessions, bit-for-bit;
  % (d) frozen-noise evaluation: encoding seed fixed at eval time (training stays stochastic);
  % (e) model selection on validation only; test set frozen.
\subsection{Triviality analysis}  % alpha-to-zero degenerate solution; "conflict reduction must be reported with retention"; per-synapse cosine degeneracy -> per-row cosine
\subsection{Ethical / boundary statement} % frozen-pretrained-component boundary (mirrors plan SS1.4)

% ------------------------------------------------------------
\section{Experimental Setup}
% RTX 5060 8GB; software version table; energy measurement convention. One environment table.

% ------------------------------------------------------------
\section{Results}
\subsection{W1: Hebbian feature learning}
  % Table: ridge readout 97.79% vs SGD readout 98.01% -- feature quality vs boundary statement
\subsection{W2: e-prop sequential learning}
  \subsubsection{Learning-rate sweep}      % Table: 4 lr x 5 seeds; 2e-3 arm reproduces September results bit-for-bit
  \subsubsection{Neuron type interaction}  % Table: LIF vs ALIF at matched lr (2e-3): 0.7867 vs 0.7764; intervals do not overlap
  \subsubsection{Budget extension}         % Figure: 250ep convergence curve -> 0.8831; plateau criterion pre-registered
  \subsubsection{Comparison with BPTT}     % \todo{headline control-arm number; matched budget; gap = X.XX pp}
\subsection{W3: R-STDP with Critic correction}
  % CartPole median 237.8 (20 seeds); with/without baseline comparison; protocol cross-ref SS3.2
\subsection{SHD: an honest negative result}
  % 5.34% vs 5% chance; attribution (hyperparams or implementation); gap vs ESPP 70.31% discussed -- as-is, not hidden
\subsection{Resource measurements}
  % memory/throughput vs n_rec curve (log-log); empirical basis for consumer-grade deployment

% ------------------------------------------------------------
\section{Discussion}
\subsection{When do local rules suffice?}    % readable conclusions by budget / task type
\subsection{The interaction matters more than the rule} % policy implication of lr interaction: hyperparameter audit should be standard practice
\subsection{Limits}                          % scale (256 -> 50k neurons unvalidated); task family (MNIST); single hardware

% ------------------------------------------------------------
\section{Conclusion}
% Three sentences: honest measurement is a prerequisite for local-learning research; this paper provides protocol + numbers + code; future: arbiter synergy benchmark.

% ------------------------------------------------------------
\bibliographystyle{plainnat}
\bibliography{refs}   % \todo{refs.bib: ~21 entries, from plan references + related_work.md}

\appendix
\section{Full configuration dumps}   % verbatim config dump per acceptance run
\section{Reproduction checklist}       % examples/paper_*.py mapped one-to-one
\section{Triviality analysis details} % alpha degeneracy sweep curves

\end{document}
